\section{Mechanism Design}
\label{sec:mech}

This section describes the inspected legacy mechanism and states only code-level or conditional
results. It deliberately does not assert that truthful voting is a strict Bayesian--Nash equilibrium,
that a $35\%$ coalition is harmless, that a Sybil attack is unprofitable, or that all funds are
exactly conserved. Those conclusions require assumptions or implementation properties absent from
the current evidence.

\subsection{Legacy mechanism}
\label{sec:mechanism}

For one disputed trade, let $m_B,m_S,m_\bot$ be the counts of BUYER, SELLER, and ABSTAIN votes.
Each address that votes deposits a fixed amount $v>0$. Define the effective-vote count
$T=m_B+m_S$; abstentions are excluded from both the threshold denominator and the quorum. The legacy
constants are
\[
  \kappa=\frac{6500}{10000}=0.65,\qquad \nu=3.
\]
Thus the mechanism does \emph{not} implement an exact two-thirds rule.

\begin{definition}[Legacy adjudication rule]
\label{def:rule}
If $T=0$, the case is void. Otherwise side $X\in\{B,S\}$ is selected when
$m_X\cdot10000\ge T\cdot6500$. A selected side produces an effective verdict only when
$T\ge\nu$; otherwise the case is void. Since $0.65>1/2$, both sides cannot satisfy the predicate
when $T>0$. In the inspected snapshot a void case triggers a full buyer-wins escrow settlement, while
all vote deposits are claimable as refunds.
\end{definition}

\begin{definition}[Legacy voter settlement]
\label{def:settlement}
In an effective case, each voter on the selected side receives
\[
  v+\left\lfloor\frac{v m_\ell}{m_w}\right\rfloor,
\]
where $m_w$ and $m_\ell$ are the winning and losing effective-vote counts. Losing voters receive
zero, and abstainers may reclaim $v$. In a void case every participant may reclaim $v$. Claims are
pull payments and therefore require each eligible address to transact successfully.
\end{definition}

\paragraph{Escrow linkage.}
The legacy escrow holds buyer amount $x$ and guarantor collateral
\[
  g=\left\lfloor\frac{x\gamma}{10^{18}}\right\rfloor.
\]
The premium $\pi$ is quoted but is not separately deposited.
The guard is $\pi\le x$, so \texttt{premium == amount} is accepted and makes the seller payment zero
on a seller-win or normal-release path; the code does not require a positive seller remainder. On a
full buyer win, the buyer is sent $x+g$. On a seller win or normal release, the seller is sent
$x-\pi$ and the guarantor is sent $g+\pi$. Partial buyer resolution sends the buyer
$\lfloor x\beta/10000\rfloor+g$ and the seller the remaining buyer principal.

\subsection{What the threshold does and does not imply}
\label{sec:analysis}

\begin{proposition}[Vote count required to select a side]
\label{prop:threshold}
Suppose $h$ effective votes oppose a coalition and all $c$ coalition votes support one side. That
side is selected by the legacy rule if and only if
\[
 c\ge\left\lceil\frac{13h}{7}\right\rceil.
\]
This is a deterministic statement about a fixed vote profile, not a profitability or truthfulness
claim.
\end{proposition}

\begin{proof}
The side is selected exactly when $c/(c+h)\ge0.65=13/20$. Rearranging gives $7c\ge13h$;
integrality gives the ceiling. The separate quorum condition must also hold.
\end{proof}

A coalition with at most $35\%$ of effective votes cannot by itself select its preferred side against
unanimous opposition. It is nevertheless not ``strategically harmless.'' If the opposing side has
less than $65\%$, neither side is selected and the case is void. Because void cases default to the
buyer, a coalition can change the economic outcome by inducing a void, particularly when it favors
the buyer. It can also bribe, split honest votes, exploit abstention, or benefit from signal errors.
The threshold provides a vote-count barrier, not general collusion resistance.

\subsection{Abstention and equilibrium status}
\label{sec:abstention}

The legacy payoff rule creates an unresolved participation problem. An abstainer receives net
monetary payoff zero after refund. A voter receives a positive bonus only if on an effective winning
side, loses $v$ if on an effective losing side, and receives zero net payoff if the case is void. If
$p_w$ and $p_\ell$ denote the probabilities that a contemplated vote is respectively winning and
losing, its expected net payoff is
\[
  p_w\left\lfloor\frac{v m_\ell}{m_w}\right\rfloor-p_\ell v,
\]
with the counts outcome-dependent. Voting beats abstention only when this expression is positive.
Signal accuracy $q>1/2$ alone does not establish that inequality, because the threshold, pivotality,
other strategies, correlated signals, and the possibility of a void all matter.

Consequently, the earlier strict-Bayesian--Nash statement and its proof are withdrawn. A future
analysis could establish a \emph{conditional} equilibrium under a fully specified prior, evidence
model, electorate-size distribution, participation reward, and strategy profile. The current score
contract records seller trade outcomes, not voter participation, so the proposed reputation loop
does not presently turn abstention into a repeated-game penalty. That integration remains future
work.

\subsection{Sybil cost is capital at risk, not proven loss}
\label{sec:sybil}

One address may vote once per case, but the inspected voting contract does not require the address to
hold a registered agent token. Creating additional voting addresses is therefore cheap. Each vote
must lock $v$ until settlement and claim, which creates a linear \emph{capital requirement}. It is
not an equal linear expenditure: a winning Sybil recovers principal and may earn a reward; a void-case
Sybil can reclaim principal; only a losing effective vote forfeits principal. Registration fees are
irrelevant to the legacy voting path unless a future release binds voters to registered identities.
Even with such a binding, whether a fee is a social cost, protocol revenue, refundable capital, or
attacker-controlled transfer depends on governance and ownership.

No ``Sybil attacks are unprofitable'' theorem follows from Proposition~\ref{prop:threshold}. Profit
also depends on the dispute value, side payments, opportunity cost and duration of locked capital,
probability of each outcome, gas, registration policy, and whether the attacker benefits from a void.
The prototype establishes only that each legacy vote requires contemporaneous stake.

\subsection{Conditional accounting}
\label{sec:accounting}

\begin{proposition}[Voting-pool claim accounting]
\label{prop:voting-accounting}
Assume a case reaches settlement; every eligible claimant claims exactly once; every transfer
succeeds; and no unrelated Ether is included. Then a void case makes claims totaling
$(m_B+m_S+m_\bot)v$. An effective case makes claims totaling
\[
 (m_B+m_S+m_\bot)v-\delta,
 \quad
 \delta=(v m_\ell)\bmod m_w,
 \quad 0\le\delta<m_w.
\]
The remainder $\delta$ remains in the contract unless a separate recovery mechanism exists.
\end{proposition}

\begin{proof}
The winning side claims
$m_wv+m_w\lfloor vm_\ell/m_w\rfloor=m_wv+vm_\ell-\delta$; abstainers claim
$m_\bot v$ and losers claim zero. The void expression follows because every participant is eligible
for principal. These are claim entitlements, not a liveness guarantee.
\end{proof}

\begin{proposition}[Escrow branch accounting]
\label{prop:escrow-accounting}
For the legacy payout formulas, if the contract enters the stated branch, all calls succeed, and its
trade-specific custody is exactly $x+g$, the nominal payout sums are $x+g$ on full buyer win,
seller win, normal release, and partial buyer resolution. Before guarantee, the funded-timeout path
pays $x$. This proposition does not prove that every reachable state has a live exit or that the
contract balance always equals the sum of trade-specific liabilities.
\end{proposition}

\begin{proof}
The sums are respectively $(x+g)$, $(x-\pi)+(g+\pi)$, the same normal-release sum, and
$\lfloor x\beta/10000\rfloor+g+x-\lfloor x\beta/10000\rfloor$. The identities permit
$\pi=x$ and hence a zero seller payment.
\end{proof}

These propositions replace ``exact conservation.'' They do not cover unclaimed balances, failed
recipient calls, forced Ether, cross-trade balance commingling, dust recovery, liveness, or future
code. Moreover, the legacy escrow performs external value transfers before setting the terminal
state in internal release helpers. Entrypoints use \texttt{nonReentrant}, but that ordering is not the
canonical checks--effects--interactions pattern and should not be described as such.

\subsection{Reputation and guarantee pricing hypothesis}
\label{sec:pricing}

The legacy reputation score is a descriptive counter-derived quantity. If
$\rho=(c,d,w,l)$ and $N=c+d+w+l>0$, the inspected code computes, with Solidity integer arithmetic,
\[
  s(\rho)=100-\left\lfloor\frac{100d}{N}\right\rfloor
             -\left\lfloor\frac{50l}{N}\right\rfloor,
\]
and returns $50$ when $N=0$. The score is not a calibrated probability and ignores transaction
value, recency, domain, buyer behavior, uncertainty, and correlated counterparties.

For illustration only, if an external underwriter independently estimates a default probability
$p_D\in[0,1)$ and loses collateral $g$ on default while earning premium $\pi$ otherwise, a zero-profit
condition under risk neutrality is
\[
 (1-p_D)\pi-p_Dg\ge0
 \quad\Longleftrightarrow\quad
 \pi\ge\frac{p_D}{1-p_D}g.
\]
This is a conditional actuarial identity, not an on-chain rule. Substituting
$p_D=(100-s)/100$ would be an unvalidated modeling choice and can imply a quote above the legacy
constraint $\pi\le x$. The current contracts neither make that substitution nor enforce the floor.
Closing the loop requires empirical calibration, uncertainty bounds, manipulation resistance, and
final end-to-end code verification.

\subsection{Threats to validity}
The analysis assumes a binary dispute despite ambiguous or partial performance; ignores evidence
quality and correlated signals; treats stake utility as linear; and omits bribery, hedging, token
volatility, gas, latency, and identity resets. Open voting reveals interim information and permits
late strategic voting. A fixed three-vote quorum is weak for high-value trades. Buyer-favoring void
settlement creates a filing and nonparticipation attack surface. Registration tokens are transferable
under ordinary ERC-721 behavior in the inspected registry, so ``soulbound'' or permanent principal
binding is not established. These limitations prevent security or welfare conclusions beyond the
narrow propositions above.
